\documentclass[10pt,a4paper]{amsart}
\usepackage[margin=19mm]{geometry}
\usepackage{amsmath,amssymb,mathtools}
\usepackage[scr=rsfs,cal=euler]{mathalfa}
\usepackage{xfrac}
\usepackage[unicode,hidelinks,bookmarks=false]{hyperref}
\newcommand{\ie}{i.e.\ }
\newcommand{\cf}{cf.\ }
\newcommand{\resp}{respectively\ }
\newcommand{\Subsection}{\subsection}
\providecommand{\nobreakdash}{\nobreak\hbox{-}}
\setlength{\emergencystretch}{2em}
\multlinegap0pt
\usepackage{bm}
\usepackage{bbm}
\usepackage{dsfont}
\usepackage{xspace}
\usepackage{cancel}
\usepackage{mathtools}
\usepackage{amsthm}
\makeatletter
\@ifundefined{thm}{\newtheorem{thm}{Thm}}{}
\@ifundefined{cor}{\newtheorem{cor}[thm]{Corollary}}{}
\@ifundefined{defn}{\newtheorem{defn}[thm]{Definition}}{}
\@ifundefined{lem}{\newtheorem{lem}[thm]{Lemma}}{}
\@ifundefined{prop}{\newtheorem{prop}[thm]{Proposition}}{}
\makeatother
\newcommand{\A}{\mathcal A}
\newcommand{\B}{\mathcal B}
\newcommand{\F}{\mathcal{F}}
\newcommand{\PA}{\mathsf{PA}}
\newcommand{\Poly}{{\hQ[{\tt x}_i]}_{i\in \hn}^{\text{def}}}
\newcommand{\ca}{\mathcal}
\newcommand{\dom}{\operatorname{dom}}
\newcommand{\ev}{\textnormal{ev}}
\newcommand{\hQ}{\mathbf Q}
\newcommand{\hn}{\mathbf N}
\newcommand{\ideal}[1]{\ca #1}
\newcommand{\img}{\operatorname{Im}}
\newcommand{\polyn}[2]{{#1}[{\tt x}_i]_{i<#2}^\text{def}}
\newcommand{\poly}[1]{{\hQ[{\tt x}_i]}_{i<#1}^\text{def}}
\newcommand{\pto}{\leadsto} %partial map
\title[Projective curves and weak second-order logic]{Projective curves and weak second-order logic}
\author{Alessandro Berarducci, Francesco Gallinaro}
\date{Source: arXiv:2503.10473v5 (CC BY 4.0)}
\begin{document}
\maketitle

\noindent\emph{Source and licence.} Projective curves and weak second-order logic. Version arXiv:2503.10473v5 (CC BY 4.0), distributed under the Creative Commons Attribution 4.0 International licence, as recorded in the arXiv licence metadata for this exact version and shown on the abstract page https://arxiv.org/abs/2503.10473v5. Full rights notice: licenses/arxiv-2503.10473-CC-BY-4.0.txt.\par\smallskip

\noindent\textit{Editorial scope.} Counted unit: the Lem whose source label is \texttt{lem:main} in the article (source0.tex lines 1130--1131, source label \texttt{lem:main}), delivered together with the article's own complete original proof of it (source0.tex lines 1133--1177, one \texttt{proof} environment of 45 lines). No mathematics is written, repaired or re-derived: statement and proof are the article's own text line for line. Required context retained uncounted: defn:piso (lines 241--249); prop:back forth (lines 253--255); prop:enumeration (lines 702--705); defn:generated-field (lines 923--933); prop:generic (lines 971--978); cor:ideal (lines 1096--1098). Every definition, display and label of those lines is retained unchanged. Omitted from this package and not redistributed: 1--240, 250--252, 256--701, 706--922, 934--970, 979--1095, 1099--1129, 1132--1132, 1178--1657. Nothing inside a retained excerpt is omitted or altered: every retained body is delivered exactly as the article's lines are written, so no omission receipt of any kind accompanies this package. Citations: the retained excerpts carry no citation command at all, so no bibliography entry of the article is transcribed and no reference is redistributed. Reference adaptation: none was needed and none was made; every reference inside the delivered unit resolves inside it, so no unresolved reference or citation is produced. Printed slips kept literal: no printed word, symbol, hypothesis or number of the counted statement or of its proof was altered. Layout and class: the article's own class is \texttt{amsart} with packages including cmbright, enumitem, amsmath, amssymb, hyperref, doi; the wrapper is the lane's pinned \texttt{amsart} editorial wrapper at 10pt on A4, which restates the theorem-like environments the retained text needs and adds the article's own macro definitions that the retained text needs (\texttt{\textbackslash A}, \texttt{\textbackslash B}, \texttt{\textbackslash F}, \texttt{\textbackslash PA}, \texttt{\textbackslash Poly}, \texttt{\textbackslash ca}, \texttt{\textbackslash dom}, \texttt{\textbackslash ev}, \texttt{\textbackslash hQ}, \texttt{\textbackslash hn}, \texttt{\textbackslash ideal}, \texttt{\textbackslash img}, \texttt{\textbackslash poly}, \texttt{\textbackslash polyn}, \texttt{\textbackslash pto}), with extra wrapper packages (bm, bbm, dsfont, xspace, cancel, mathtools). The delivered numbering is the unit's own. Assets: no retained span contains a figure environment, a graphics-inclusion command, a PDF-page inclusion, a table or a TikZ picture; no image file is copied into this package, the package declares no asset dependency and no image file is redistributed with it. Licence: version arXiv:2503.10473v5 of this article is distributed under the Creative Commons Attribution 4.0 International licence, as recorded in the arXiv licence metadata for this exact version and shown on the abstract page https://arxiv.org/abs/2503.10473v5. A full rights notice is delivered at licenses/arxiv-2503.10473-CC-BY-4.0.txt. Characters: every retained excerpt is pure ASCII, so the delivered engine, XeLaTeX with its default Latin Modern text font, reports no missing character.
\begin{defn}\label{defn:piso}
We say that a partial map $\iota: \ca M \pto \ca N $ between two $L$-structures is a {\em partial isomorphism} if for every $a_1, \ldots, a_n \in \dom(\iota )$ and every quantifier free formula 
%
$\varphi(x_1, \ldots, x_n)$, with variables appropriate for the sorts of $a_1, \ldots, a_n$, we have: 
%
\begin{equation*}
\ca M \models \varphi(a_1, \ldots, a_n) \iff \ca N \models \varphi (\iota (a_1), \ldots, \iota (a_n))
\end{equation*}
\end{defn}

\begin{prop}\label{prop:back forth}
If $\ca G$ is a family of partial isomorphisms $\iota: \ca M \pto \ca N$ with the back and forth property, then every $\iota \in \ca G$ is an \emph{elementary map}, namely the equivalence in Definition \ref{defn:piso} holds for every formula of the language, not only for the quantifier free formulas. If such a family $\ca G$ exists and is non-empty, then $\ca M$ and $\ca N$ are elementarily equivalent. 
\end{prop}

\begin{prop}\label{prop:enumeration}
Let $\ca N \models \PA(\ca A)$. 
If $X$ is a hyperfinite set in $\ca A$ and $|\ca N| \times X$ has a definable finite power set in $\ca A$, then there exist a unique $n\in |\ca N|$ and a hyperfinite bijection $f:(n)\to X$.
\end{prop}

\begin{defn}\label{defn:generated-field}
Let $n\in \hn$ and let $f\in K^{(n)}$. We define 
$$\hQ[f]^{\text{def}} = \{\ev(p, f) \mid p\in \polyn {\hQ} n\} \subseteq K$$ 
and we let $\hQ(f)^{\text{def}} \subseteq K$ be the quotient field of the ring $\hQ[f]^{\text{def}}$. 

It is easy to see that $\hQ[f]^{\text{def}}$ depends only on the image of $f$. Moreover, for every $a\in \F(K)$ there is $n\in \hn$ and $f\in K^{(n)}$ such that $a = \img(f)$, so we can define 
$$\hQ[a]^{\text{def}} = \hQ[f]^{\text{def}}.$$ We call $\hQ[a]^{\text{def}}$ the definable subring generated by $a$ and we call $\hQ(a)^{\text{def}} = \hQ(f)^{\text{def}}$ the definable subfield generated by $a$. 

In a similar way we define the definable subring $\hQ[S]^{\text{def}} \subseteq K$ generated by a finite subset $S = \{a_1, \ldots, a_n, b_1, \ldots, b_m\}$ of $K\cup \F(K)$ where $a_1, \ldots, a_n\in K$ and $b_1, \ldots, b_m\in \F(K)$ (with $n,m$ standard): it suffices to let $A = \{a_1\}^* \cup^* \ldots, \cup^* \{a_n\}^* \cup^* b_1 \cup^* \ldots \cup^* b_m \in \F(K)$
and put $\hQ[S]^{\text{def}} = \hQ[A]^{\text{def}}$. Similarly we define $\hQ(S)^{\text{def}} = \hQ(A)^{\text{def}}$. 
\end{defn}

\begin{prop}[Existence of generic zeros] \label{prop:generic} Let $R$ be a definable subfield of $K$ generated by a hyperfinite set (Definition \ref{defn:generated-field}). Suppose that $\ca K = (K, \F(K))$ is definably algebraically closed and of hyper-infinite transcendence degree. 
Let $n\in \hn$ and let $\ideal I \subseteq \polyn R n$ be a definable prime ideal. Then:  
\begin{enumerate}
\item 
$\ideal I$ has a generic zero. 
\item Let $r\leq n$ and let ${\ideal I}_r = \ideal I \cap \polyn R r$. Then every generic zero $\bar \alpha \in K^{(r)}$ of ${\ideal I}_r$ can be extended to a generic zero $\bar \alpha \bar \beta \in K^{(n)}$ of $\ideal I$. 
\end{enumerate}
\end{prop}

\begin{cor} \label{cor:ideal}  
Every ideal $\ideal I \subseteq \poly n$ definable in $(K, \F(K))^{eq}$ is definable in $(\hQ, \hn,+,\cdot)^{eq}$. 
\end{cor} 

\begin{lem} \label{lem:main} Let $\mathcal A= (A, \F (A))$ and $\mathcal B = (B, \F(B))$ be models of $T_{\text{rec}}$. Suppose that $\Psi: \hn(\ca A) \cong \hn (\ca B)$ is an isomorphism in the language $\{+,\cdot\}$. Then $\Psi$ is an elementary map when considered as a partial function from $\ca A$ to $\ca B$. In particular $\ca A \equiv \ca B$. 
\end{lem} 

\begin{proof} It suffices to show that the restriction of $\Psi$ to any finite subset of $\hn (\ca A)$ belongs to a family of partial isomorphisms $\ca G$ from $\ca A$ to $\ca B$
 with the back and forth property. Replacing $\ca A$ with an isomorphic structure, we may assume that $\hn(\ca A) = \hn(\ca B)$, and $\Psi$ is the identity on $(\hn,+,\cdot) = (\hn (\ca A),+,\cdot) = (\hn (\ca B),+,\cdot)$. 
We can then assume that $\hQ(\ca A) = \hQ(\ca B)$ and 
write
%
$$(\hQ, \hn, +, \cdot)^{eq} = (\hQ(\ca A), \hn (\A), +,\cdot)^{eq} = (\hQ(\ca B), \hn (\B), +,\cdot)^{eq}$$
%
In particular, since $\Poly$ and its operations are definable in $(\hQ, \hn, +,\cdot)^{eq}$, we may assume that $\ca A$ and $\ca B$ have the same non-standard polynomials over $\hQ$. 

The idea is to show that the relations between elements of $A\cup \F(A)$ can be coded by definable ideals in $\poly n$ (for various $n$) and such ideals are definable in $(\hQ,\hn, +, \cdot)^{eq}$ (Corollary \ref{cor:ideal}). We will use such ideals to define a family $\ca G$ of partial isomorphisms with the back and forth property. 
To this aim we observe that an element $a \in \F(A)$ can be coded by a function $f\in A^{(n)}$ with $a = \img(f)$ (Proposition \ref{prop:enumeration}) where $n\in \hn$ can be non-standard. Thus we need to ensure that the relations between elements of   $\bigcup_{n \in \hn} A^{(n)}$ be preserved by the partial maps in the family. This motivates the following definition. 

Let $\ca G_0$ be the set of all partial maps $\iota : \bigcup_{n \in \hn} A^{(n)} \pto \bigcup_{n\in \hn} B^{(n)}$ with the following properties: 
\begin{itemize}
\item $\dom(\iota)$ is finite (in the standard sense). 
\item If $n\in \hn$ and $f\in A^{(n)} \cap \dom(\iota)$ and then $\iota(f) \in B^{(n)}$. 
\item If $\{f_1, \ldots, f_k\} \subseteq \dom(\iota)$ and $f = f_1*\ldots *f_k \in A^{(m)}$ (where ``$*$'' is the concatenation), the definable ideal ${\ideal I}_f \subseteq \poly m$ of $f$ in $\ca A$ coincides with the definable ideal ${\ideal I}_g$ of $g = \iota(f_1)*\ldots * \iota(f_k)$ in $\ca B$. 
\end{itemize}

We claim that $\ca G_0$ has the back and forth property. 
Given $\iota$ as above with $\dom(\iota) = \{f_1, \ldots, f_k\}$ and $\img(\iota) = \{g_1, \ldots, g_k\}$, consider a new $f$ and we need to find a corresponding $g$. Let $\ideal I$ be the definable ideal of  $f_1*\ldots f_k*f$ in $\ca A$. Then $\ideal I$ is definable in $(\hQ,\hn, +,\cdot)^{eq}$ (Corollary \ref{cor:ideal}), hence it is also definable in $\ca B$ (by the same formula) and therefore it has a generic zero of the form $g_1* \ldots * g_k*g$ (Proposition \ref{prop:generic}(2)). This proves the forth direction and the back direction is analogous. 

\bigskip 
Given $\iota \in \ca G_0$, let $\hat \iota: A \cup \F(A) \pto B \cup \F(B)$ be the smallest function such that, for all $f\in A^{(n)}$ and $m<n$, if $\iota(f) = g$, then $\hat \iota$ maps $f(m)\in A$ to $g(m)\in B$ and $\img (f) \in \F(A)$ to $\img(g) \in \F(B)$. 
To verify that $\hat \iota$ is well defined, suppose that $f_1\in A^{(n_1)}$ and $f_2\in A^{(n_2)}$ belong to $\dom(\iota)$ and $f_1(m_1) = f_2(m_2)$ (where $f_1,f_2$ are not necessarily distinct). Then the definable ideal of $f_1*f_2$ contains the polynomial $\tt x_{m_1} - \tt x_{n_1+m_2}$. This polynomial then lies in the definable ideal of $g_1*g_2$ where $g_1 = \iota(f_1)$ and $g_2 = \iota(f_2)$. Since $\iota$ preserves the information contained in the ideal, $g_1(m_1) = g_2(m_2)$. Similarly one shows that if $\img(f_1) = \img(f_2)$, then $\img(g_1) = \img(g_2)$. The same arguments show that $\hat \iota$ is injective. Let us also observe that if $x\in \hQ$ belongs to the domain of $\hat \iota$, then $\hat \iota (x) = x$. 


Now let $\ca G = \{\hat \iota \mid \iota \in \ca G_0\}$. Then $\ca G$ is a family of injective partial maps from $A\cup \F(A)$ to $B \cup \F(B)$ and it has the back and forth property because $\ca G_0$ does. 

We claim that every element $\hat \iota$ of $\ca G$ is a partial isomorphism, namely it preserves the quantifier free formulas, or equivalently the atomic formulas (negations can be handled by considering the inverse of $\hat \iota$). So we must prove: 
%
$$
\begin{cases}
x\in a \implies \hat \iota(x) \in \hat \iota(a)\\
x+y = z \implies \hat \iota(x) + \hat \iota(y) = \hat \iota(z)\\
x\cdot y= z \implies \hat \iota(x)\cdot \hat \iota(y) = \hat \iota(z)
\end{cases}
$$
Suppose for instance that $a\in \dom(\hat \iota)$ and $x\in a$. Then there is $f\in \dom(\iota)$ such that $a = \img(f)$ and there is $m\in \hn$ such that $x = f(m)$. Let $g = \iota(f)$. Then by definition of $\hat \iota$ we have $\hat \iota (x) = g(m) \in \img (g) = \hat \iota (\img f)) = \hat \iota (a)$. 

Now suppose that $x,y,z\in \dom(\hat \iota)$ and $x+y= z$. Then we can write $x = f_1(m_1), y=f_2(m_2), z= f_3(m_3)$ where $f_1 \in A^{(n_1)}$, $f_2\in A^{(n_2)}$, $f_3 \in A^{(n_3)}$ are in $ \dom(\iota)$ (not necessarily distinct) and $m_i<n_i$ for $i=1,2,3$. Let $g_i = \iota(f_i)$. Then the polynomial $\tt x_{m_1} + \tt x_{n_1+m_2} - \tt x_{n_1+n_2+m_3}$ belongs to the definable ideal of $f_1*f_2*f_3$, hence also to the definable ideal of $g_1*g_2*g_3$. It follows that $\hat \iota (x) + \hat \iota(y) = g_1(m_1) + g_2(m_2) = g_3(m_3) =  \hat \iota (z)$. 
The proof of $x\cdot y= z \implies \hat \iota(x)\cdot \hat \iota(y) = \hat \iota(z)$ is analogous. 

Since for any $n\in \hn \cap \dom(\hat \iota)$, we have $\hat \iota (n) = n = \Psi(n)$, it follows that every finite restriction of $\Psi$ can be extended to a map in $\ca G$ and therefore $\Psi$ is an elementary map by Proposition \ref{prop:back forth}. 
\end{proof}

\end{document}
